%! Describing the Distribution of a Quantitative Variable — Unit 1, Lesson 1.4 — Slides (Beamer)
\documentclass[aspectratio=169, 11pt]{beamer}
\usepackage{apstats-beamer}   % provides \navyheader{} and \sectionlabel[color]{}
\usetikzlibrary{positioning}  % -beamer loads tikz but not this library

% The C and Y column types live in apstats-article, which a beamer deck does not
% load — redeclare them here rather than pulling in the article preamble.
\newcolumntype{C}[1]{>{\centering\arraybackslash}p{#1}}
\newcolumntype{Y}{>{\centering\arraybackslash}X}

\begin{document}

% TITLE SLIDE (CourseName is not defined in beamer, so the name is literal)
{
\setbeamercolor{background canvas}{bg=navy}
\begin{frame}[plain]
  \vspace{2.8cm}\hspace{0.55cm}%
  \begin{minipage}{0.9\textwidth}
    {\color{goldacc}\bfseries\small LESSON 1.4}\\[0.5cm]
    {\color{white}\bfseries\LARGE Describing the Distribution of a Quantitative Variable}\\[1.2cm]
    {\color{skymid}\small Statistics~$\cdot$~Unit 1}
  \end{minipage}
\end{frame}
}

% GRADUAL RELEASE deck: title → targets → warm-up hook → I DO (notes frames) → WE DO (guided
% practice) → spoken debrief → close & START THE HOMEWORK, which is the individual
% practice — there is no independent practice set in the notes → AP Practice → close.
% There is NO group activity in this course — the notes carry the whole release. There is no
% transfer set and no reflection box: the commute record lives only in the AP Practice page.
% Vocabulary is given up front here, not withheld: the notes frames ARE the direct instruction.
\newcommand{\redink}[1]{{\color{redacc}\bfseries #1}}

% ── LEARNING TARGETS ─────────────────────────────────────────────────────────
\begin{frame}[t, plain]
  \navyheader{Today's targets}
  \sectionlabel{Lesson 1.4}
  By the end of today, I can\ldots
  \begin{itemize}\setlength{\itemsep}{5pt}
    \item explain why one number is never a description of a set of data.
    \item name a distribution's \textbf{shape} --- symmetric, skewed left, or skewed right.
    \item point out the \textbf{unusual features}: outliers, gaps, and clusters.
    \item describe a whole distribution in \textbf{four parts}, in context.
  \end{itemize}
  \begin{block}{How today works}
    Warm-up $\to$ \textbf{notes together} $\to$ \textbf{one we work as a class} $\to$
    debrief $\to$ \textbf{you start the homework here, alone}.
  \end{block}
\end{frame}

% ── HOOK ─────────────────────────────────────────────────────────────────────
{
\setbeamercolor{background canvas}{bg=navy}
\begin{frame}[plain]
  \vspace{1.5cm}
  \begin{center}
    {\color{skymid}\bfseries\small TWO CHECKOUT LINES. THE SAME MIDDLE WAIT.}\\[0.9cm]
    {\color{white}\bfseries\Large One of them you would never stand in again.}\\[0.9cm]
    {\color{goldacc}\small Both managers report ``about 4 minutes.'' Only one of them is
    telling you anything.}
  \end{center}
\end{frame}
}

% ── I DO — direct instruction begins ─────────────────────────────────────────
{
\setbeamercolor{background canvas}{bg=navy}
\begin{frame}[plain]
  \vspace{1.8cm}\hspace{0.8cm}%
  \begin{minipage}{0.86\textwidth}
    {\color{goldacc}\bfseries\small GUIDED NOTES \ \textbullet\ \ I DO}\\[0.7cm]
    {\color{white}\bfseries\Large Open to your Guided Notes page. Fill each blank as we
    name it --- not before.}\\[0.9cm]
    {\color{skymid}\small Five rows go in the vocabulary box today. Write each term the
    moment we use it.}
  \end{minipage}
\end{frame}
}

% ── SECTION 1: CENTER AND VARIABILITY ────────────────────────────────────────
\begin{frame}[t, plain]
  \navyheader{Same center. Nothing else the same.}
  \sectionlabel{notes, Center \& Variability $\cdot$ problems 1--4}
  \vspace{4pt}
  \begin{columns}[t]
    \begin{column}{0.47\textwidth}
      \begin{block}{Center}
        The middle, typical value --- where the data pile up (EK UNC-1.H.5).
      \end{block}
      \vspace{4pt}
      {\small Line A: \textbf{4 minutes.} \quad Line B: \textbf{4 minutes.}}
    \end{column}
    \begin{column}{0.47\textwidth}
      \begin{block}{Variability (spread)}
        How far the values reach from each other (EK UNC-1.H.6).
      \end{block}
      \vspace{4pt}
      {\small Line A: every wait \textbf{3 to 6}. \quad Line B: \textbf{1 to 13}.}
    \end{column}
  \end{columns}
  \vspace{12pt}
  \begin{center}
    {\large ``The typical wait is 4 minutes in both lines'' is \redink{true} --- and useless.}\\[6pt]
    {\small\color{slate} A center reported without the spread beside it is a summary of nothing.}
  \end{center}
\end{frame}

% ── SECTION 1: SHAPE — PEAKS AND THE MIRROR ──────────────────────────────────
\begin{frame}[t, plain]
  \navyheader{Shape: count the peaks, then check the mirror}
  \sectionlabel{notes, Peaks \& Mirror $\cdot$ three pictures, three stories $\cdot$ problems 5--7}
  \vspace{6pt}
  \begin{columns}[t]
    \begin{column}{0.30\textwidth}
      \begin{block}{Unimodal}
        One main peak.
      \end{block}
    \end{column}
    \begin{column}{0.30\textwidth}
      \begin{block}{Bimodal}
        Two clear peaks.
      \end{block}
    \end{column}
    \begin{column}{0.30\textwidth}
      \begin{block}{Uniform}
        No peak at all.
      \end{block}
    \end{column}
  \end{columns}
  \vspace{14pt}
  \begin{center}
    {\large \textbf{Symmetric:} the left half is roughly a \redink{mirror image} of the
    right.}\\[8pt]
    {\small\color{slate} Home prices lean. Quiz scores lean the other way. Club heights do
    not lean at all. Fill the table: where does it pile, which way do the stragglers run?}
  \end{center}
\end{frame}

% ── SECTION 2: WHAT SITS APART ───────────────────────────────────────────────
\begin{frame}[t, plain]
  \navyheader{What sits apart}
  \sectionlabel{notes, Unusual Features $\cdot$ outliers, gaps, clusters $\cdot$ problems 8--11}
  \vspace{4pt}
  \begin{itemize}\setlength{\itemsep}{7pt}
    \item \textbf{\color{navy}Outlier} --- a value unusually far from the rest of the data.
          Line B's customers at \textbf{11} and \textbf{13} minutes.
    \item \textbf{\color{navy}Gap} --- a stretch of the number line where nothing was observed.
          Line B: everything from about 7 minutes to 11.
    \item \textbf{\color{navy}Cluster} --- values bunched together, set off from the others by
          a gap.
  \end{itemize}
  \vspace{10pt}
  \begin{block}{Four things or nothing: SOCS}
    \textbf{S}hape $\cdot$ \textbf{O}utliers and other unusual features $\cdot$
    \textbf{C}enter $\cdot$ \textbf{S}pread --- every one of the four \redink{in context},
    naming the variable and its units.
  \end{block}
\end{frame}

% ── THE CRUX ─────────────────────────────────────────────────────────────────
\begin{frame}[t, plain]
  \navyheader{Skewed which way? The tail decides.}
  \sectionlabel{notes, The Tail $\cdot$ the whole point of today $\cdot$ problems 12--13}
  \vspace{6pt}
  \begin{itemize}\setlength{\itemsep}{8pt}
    \item Picture I: $6+14+10 = 30$ of the 40 homes sold below \$400{,}000.
    \item And the picture still stretches past \$700{,}000 --- because two houses really did.
    \item The pile is on the \textbf{left}. The long thin tail runs to the \textbf{right}.
    \item So Picture I is \redink{skewed right.}
  \end{itemize}
  \vspace{10pt}
  \begin{block}{Essential Knowledge (UNC-1.H.3)}
    \textbf{A skew is named for the side the tail runs out on} --- not the side the pile
    sits on.
  \end{block}
\end{frame}

\begin{frame}[t, plain]
  \navyheader{Two more things before you call it done}
  \sectionlabel{notes, The Tail $\cdot$ problems 14--15, the counterweights --- do not skip these}
  \vspace{8pt}
  \begin{columns}[t]
    \begin{column}{0.46\textwidth}
      \begin{block}{Symmetric is not the whole answer}
        Picture III mirrors --- and it is also \textbf{bimodal}, with nobody at all between
        66 and 70 inches. ``Typical height about 68 inches'' describes \redink{no one} in the
        club.
      \end{block}
    \end{column}
    \begin{column}{0.46\textwidth}
      \begin{block}{An outlier is a finding, not an error}
        Line B's 13-minute customer: check her, do not \redink{delete} her. A real customer
        with a full cart is exactly what the manager needs to hear about.
      \end{block}
    \end{column}
  \end{columns}
  \vspace{14pt}
  \begin{center}
    {\large Now say all four for Picture I --- in context, with units.}
  \end{center}
\end{frame}

% ── WE DO ────────────────────────────────────────────────────────────────────
\begin{frame}[t, plain]
  \navyheader{Guided Practice: The Sleep Survey}
  \sectionlabel[goldacc]{We do $\cdot$ you hold the pen}
  Thirty students reported their hours of sleep on a school night.
  \begin{itemize}\setlength{\itemsep}{5pt}
    \item \textbf{16.} Which way does it lean? Point at the tail, then name it.
    \item \textbf{17.} How many students? How many peaks? Which value sits apart, and where
          does the gap run?
    \item \textbf{18.} What does that longer tail mean \emph{for these students}?
    \item \textbf{19.} Now say all four: shape, unusual features, center, spread. In context.
  \end{itemize}
  \begin{block}{The question behind the question}
    The pile is on the right and the tail runs left. If you answered \redink{skewed right},
    you read the pile.
  \end{block}
\end{frame}

% ── YOU DO — the homework is the individual practice, started in class ───────
\begin{frame}[t, plain]
  \navyheader{You do: start the homework}
  \sectionlabel{Last 8 minutes $\cdot$ on your own $\cdot$ finish it at home}
  \begin{itemize}\setlength{\itemsep}{7pt}
    \item Four displays, four shape descriptions --- match them.
    \item A classmate calls Display B ``skewed left, because the data is piled on the left.''
          Correct them.
    \item A two-humped picture of ages at a swim session --- name the feature, explain it.
    \item Twenty-four cereals' sugar content: write the whole four-part description.
  \end{itemize}
  \begin{block}{Silent and alone --- I am circulating}
    \emph{Item 2} is the one I am reading over your shoulder. If you agreed with the classmate,
    flag me before you leave.
  \end{block}
\end{frame}

% ── DEBRIEF ──────────────────────────────────────────────────────────────────
\begin{frame}[t, plain]
  \navyheader{Debrief: point at the tail}
  \sectionlabel{8 minutes $\cdot$ pens down, papers up --- we say these aloud}
  \vspace{6pt}
  \begin{columns}[t]
    \begin{column}{0.46\textwidth}
      \begin{block}{Picture I --- skewed right}
        Pile on the left. Tail to the right. The tail wins the name.
      \end{block}
    \end{column}
    \begin{column}{0.46\textwidth}
      \begin{block}{Picture II --- skewed left}
        Pile on the right. Tail to the left. Same rule, mirrored.
      \end{block}
    \end{column}
  \end{columns}
  \vspace{14pt}
  \begin{block}{Two things to say out loud}
    Describe Picture I in one sentence a reader could act on. Then say what those 40 homes
    \redink{cannot} tell you --- nothing about the next town, and nothing about next year.
  \end{block}
\end{frame}

% ── HOMEWORK (started in class, scored) ─────────────────────────────────────
\begin{frame}[t, plain]
  \navyheader{Homework --- start it now}
  \sectionlabel[goldacc]{5 exercises $\cdot$ scored $\cdot$ due the first class after two study halls}
  \begin{itemize}\setlength{\itemsep}{5pt}
    \item Match four displays to four descriptions.
    \item Correct the classmate who read the pile instead of the tail.
    \item Name the unusual feature in a two-humped age distribution.
    \item Describe 24 cereals' sugar content --- all four parts, in context.
    \item One multiple-choice item, with a one-sentence justification.
  \end{itemize}
  \begin{block}{This one is graded}
    This is your \redink{scored} homework, on paper, due the first class after two study halls.
    Start with 1, 2, and 5 while I am still here to answer questions --- finish the rest at
    home.
  \end{block}
\end{frame}

% ── AP PRACTICE & PREVIEW (assigned, unscored) ──────────────────────────────
\begin{frame}[t, plain]
  \navyheader{AP Practice \& what's next}
  \sectionlabel[goldacc]{AP Practice --- assigned, not scored}
  \begin{itemize}\setlength{\itemsep}{5pt}
    \item Four multiple-choice items on a 40-employee commute record.
    \item One free-response set: classify, describe, take apart a manager's one number, and
          regroup at a wider interval.
    \item Watch distractor (B) on item 2. It is today's error, written out in full.
  \end{itemize}
  \begin{block}{Not graded --- but this is what the exam looks like}
    The AP Practice page is \textbf{not scored}. Work it for the reps, and bring what you
    get stuck on to study hall or office hours. The \redink{graded} homework is the section
    behind it.
  \end{block}
\end{frame}

% ── CLOSE ────────────────────────────────────────────────────────────────────
{
\setbeamercolor{background canvas}{bg=navy}
\begin{frame}[plain]
  \vspace{1.5cm}\hspace{0.8cm}%
  \begin{minipage}{0.86\textwidth}
    {\color{goldacc}\bfseries\small BEFORE YOU GO}\\[0.7cm]
    {\color{white}\bfseries\Large A skew is named for its tail. A center without a spread is a
    summary of nothing. Say all four, in context.}\\[0.9cm]
    {\color{skymid}\small Next: Lesson 1.5, Summary Statistics --- ``about 8 hours'' and ``most
    between 6 and 9'' become the mean, the median, the IQR, and the standard deviation, and we
    find out which of them an outlier can drag around.}
  \end{minipage}
\end{frame}
}

\end{document}
